\documentclass[11pt]{article}
\usepackage[margin=1in]{geometry}
\usepackage{amsmath,amssymb,amsthm,graphicx}
\usepackage[colorlinks=true,linkcolor=blue,urlcolor=blue,citecolor=blue]{hyperref}
\hypersetup{pdftitle={One edge and a fractional graph Laplacian},
pdfauthor={Research working manuscript}}
\newtheorem{theorem}{Theorem}[section]
\newtheorem{lemma}[theorem]{Lemma}
\newtheorem{corollary}[theorem]{Corollary}
\newtheorem{proposition}[theorem]{Proposition}
\newtheorem{remark}[theorem]{Remark}
\newcommand{\R}{\mathbb R}
\newcommand{\norm}[1]{\left\lVert#1\right\rVert}
\newcommand{\tr}{\operatorname{tr}}
\newcommand{\rank}{\operatorname{rank}}
\newcommand{\diag}{\operatorname{diag}}
\newcommand{\spec}{\operatorname{spec}}
\newcommand{\N}{\mathcal N}
\newcommand{\one}{\boldsymbol 1}
\renewcommand{\labelitemi}{$\bullet$}
\title{One edge and a fractional graph Laplacian:\\
relative spectra, degree constraints, and constructive compression}
\author{Research working manuscript}
\date{November 15, 2026}
\begin{document}
\maketitle
\begin{abstract}
How many relative spectral directions can change when one graph edge is
inserted or its conductance is increased? For $0<\alpha<1$ we study
$C=(L^\dagger)^{\alpha/2}(L'^\alpha-L^\alpha)(L^\dagger)^{\alpha/2}$.
The edge leverage $\tau$ fixes the exact budget
$\sum_j\log(1+\lambda_j(C))=\alpha\log(1+\tau)$.
Every finite nonnegative spectral list exhausting this budget is a limit
on weighted paths with unit vertex masses. The same universality at
$\tau=1$ holds for one actual edge insertion between simple unweighted
connected bipartite graphs, with resistance exactly one in every host.
This gives the exact unrestricted worst-case rank above any strict
tolerance. With growing dimension, the complete finite-target uniform
closure permits a weak budget inequality; this differs from fixed-dimensional
matrix closure and from exact finite attainment.
Among simple unweighted hosts, for doubling one existing unit edge,
maximum degree two has sharp worst-case rank
$\log^2(1/\varepsilon)/(2\pi^2)$ to leading order, whereas every fixed
maximum degree $d\ge3$ has sharp leading rank
$\alpha\log(2d/(d+2))/\varepsilon$, even on bipartite bridge examples.
The upper bound subtracts the entire rooted local response using a
free/wired resistance comparison. Local constraints prevent full
spectral universality at fixed degree.
Finally, an explicit factor can be constructed with two-sided relative
error in expected nearly linear arithmetic work. All solves, feasible
flows, small-Gram computation, and factor entries are charged; the
hidden logarithms contain no graph spectral gap or conductance ratio.
For this simple unweighted existing-edge doubling problem at fixed
degree $d\ge2$, the construction attains the sharp uniform leading
rank constant. The arithmetic model and the separately charged
baseline inverse are specified. Classical matrix inequalities, rational approximation,
integral-operator asymptotics, graph coarse graining, and electrical-flow
solvers are credited throughout. Historical priority is not certified.
\end{abstract}

\section{The single-edge graph problem}\label{sec:introduction}
Let $G$ be a finite connected weighted undirected graph on $n$ vertices,
with ordinary unit vertex masses and combinatorial Laplacian $L$.
Insert an edge or increase one conductance:
\[
 L'=L+d b b^T,\qquad b=e_p-e_q,\qquad d>0.
\]
For a fixed exponent $0<\alpha<1$, define
\begin{equation}\label{eq:graph-response}
 \Delta=L'^\alpha-L^\alpha,\qquad
 C=(L^\dagger)^{\alpha/2}\Delta(L^\dagger)^{\alpha/2},\qquad
 \tau=d b^TL^\dagger b.
\end{equation}
The parameter $\tau$ is the increment times the effective resistance
between the edited endpoints in the host graph. Every inverse fractional
power in this paper is zero on constants and has the stated negative
power on $\mathbf1^\perp$. Determinants of singular graph matrices are
always taken on that common positive support. Norms are Euclidean
operator norms unless indicated otherwise. Write
\[
 \mathcal N_C(\varepsilon)=\#\{j:\lambda_j(C)>\varepsilon\},\qquad
 M_\tau(\varepsilon)=\max\left\{0,
 \left\lceil\frac{\alpha\log(1+\tau)}{\log(1+\varepsilon)}\right\rceil-1
 \right\}.
\]
Strict inequalities are used throughout, including at integral thresholds.
For graph pairs, we also write $\mathcal C_\alpha(L,L')=C$.

This metric measures the correction against the original fractional
energy. If $C_r$ is its positive spectral truncation and
$F=L^{\alpha/2}C_rL^{\alpha/2}$, then
\begin{equation}\label{eq:one-sided-goal}
 0\preceq\Delta-F\preceq\varepsilon L^\alpha
 \quad\Longleftrightarrow\quad
 \norm{C-C_r}\le\varepsilon.
\end{equation}
Consequently $\mathcal N_C(\varepsilon)$ is exactly the minimum rank
for this one-sided structural approximation. It is also the minimum
rank for unrestricted operator-norm approximation of $C$ to tolerance
$\varepsilon$. Absolute error for $\Delta$, or error divided by
$\norm\Delta$, is a different criterion.

\begin{proposition}[Exact budget and universal rank bound]\label{prop:budget}
For every positive single-edge edit,
\[
 C\succeq0,\qquad
 \sum_j\log(1+\lambda_j(C))=\alpha\log(1+\tau),\qquad
 \mathcal N_C(\varepsilon)\le M_\tau(\varepsilon).
\]
For each fixed $\tau>0$ the last bound is attained as a worst-case
strict count on unit-mass weighted paths, and
$M_\tau(\varepsilon)=\alpha\log(1+\tau)/\varepsilon+O_{\alpha,\tau}(1)$.
\end{proposition}
\begin{proof}
Restrict to $\mathbf1^\perp$ and set $A=L$, $u=\sqrt d\,b$.
The fractional resolvent formula proves positivity. The determinant
lemma gives
\[
 \det(A^{-\alpha/2}(A+uu^T)^\alpha A^{-\alpha/2})
 =(1+u^TA^{-1}u)^\alpha.
\]
If $k$ correction eigenvalues exceed $\varepsilon$, their logarithms
give $k\log(1+\varepsilon)<\alpha\log(1+\tau)$.
The sharpness is proved in Theorem~\ref{w:thm:path} and
Corollary~\ref{mat:cor:count} below.
\end{proof}

The graph class changes what can occur within this budget. The main
results, all for unit vertex masses, are as follows.
\begin{center}
\small
\begin{tabular}{@{}p{0.31\textwidth}p{0.32\textwidth}p{0.29\textwidth}@{}}
\hline
Host and edit & Spectral conclusion & Worst-case strict count\\
\hline
Weighted path; one bridge increased by factor $1+\tau$
& Every full-budget finite list is a limit
& Exactly $M_\tau(\varepsilon)$\\[3pt]
Simple unweighted bipartite host and update; one absent edge inserted
with resistance exactly one
& Every list of product $2^\alpha$ is a limit
& Exactly $M_1(\varepsilon)$\\[3pt]
Simple unweighted path or cycle; one existing edge doubled
& Uniform leverage-measure control
& $\sim\log^2(1/\varepsilon)/(2\pi^2)$\\[3pt]
Simple unweighted bipartite host of maximum degree three; one bridge doubled
& Sharp local/global budget
& $\sim\alpha\log(6/5)/\varepsilon$\\
\hline
\end{tabular}
\end{center}
The first two rows concern limits, not arbitrary exact finite
realization. The second permits growing graph size and degree. The last
two concern an existing edge doubled to conductance two, so the updated
graph has one weighted edge. Their combinatorial degree is unchanged.
For every fixed $d\ge3$, allowing any existing edge to be doubled gives
\[
 \lim_{\varepsilon\downarrow0}\varepsilon W_{d,\alpha}(\varepsilon)
   =A_{d,\alpha}:=\alpha\log\frac{2d}{d+2}.
\]
The same coefficient is attained in the bipartite bridge subclass.
This is an exact leading asymptotic, rather than a finite-threshold count
formula. At fixed degree, local response prevents full spectral universality.
For degree two, the sharp coefficient is $1/(2\pi^2)$, independent of
the fixed $\alpha$. Thus the leading worst-case rank is determined
at every fixed maximum degree $d\ge2$.

\begin{figure}[tbp]
\centering
\includegraphics[width=.92\textwidth]{figures/relative-rank-degree-transition.pdf}
\caption{The degree dichotomy for doubling an existing unit edge. Reservoir
circles represent large finite bipartite graphs, with their internal
vertices suppressed. Rank means the strict energy-relative count;
the degree laws are proved in Section~\ref{deg:sec:degree}. The sharp
coefficient uses additional repeated blocks and multiple reservoir ports,
which are not depicted in this schematic.}
\label{fig:degree}
\end{figure}

There are three distinct spectral assertions. In fixed matrix dimension,
the rank-one closure is exactly
$q_i\ge1$, $\prod_iq_i=(1+\tau)^\alpha$.
For a fixed finite target padded by ones while graph dimension grows,
uniform spectral closure instead requires
$\prod_iq_i\le(1+\tau)^\alpha$; arbitrarily many individually tiny
entries can carry the missing budget. On the full-budget face the
constructions can be chosen to converge in $\ell^1$ after subtracting
one. Exact finite matrix attainment imposes another restriction:
the entries above one must be pairwise distinct
(Appendix~\ref{ex:sec:finite-attainment}).

The constructive counterpart is Theorem~\ref{alg:thm:main}. Given the
explicit weighted graph and a positive edit, it returns an explicit
$X\in\R^{n\times r}$ with
\[
 -\frac\varepsilon8L^\alpha\preceq\Delta-XX^T
       \preceq\frac{3\varepsilon}4L^\alpha,
 \qquad
 r<\frac{\alpha\log(1+\tau)}{\log(1+\varepsilon/4)}\quad(r>0).
\]
Its expected arithmetic work is $\widetilde O_\alpha(m+n)$, with
logarithms only in $n,1+\tau,1/\varepsilon,1/\delta$.
This computed approximation is two-sided. It is not asserted to be
the one-sided optimal truncation in~\eqref{eq:one-sided-goal}.
The algorithm reconstructs from the current graph per edit; it is not
a sublinear maintained data structure. Baseline inverse fractional
actions for applying the resulting preconditioner are charged separately.
For simple unweighted hosts with one existing edge doubled at fixed
maximum degree $d\ge2$,
Corollary~\ref{alg:cor:sharp-degree} attains the corresponding sharp
leading rank in both regimes, uniformly over the graph class,
with relative operator error at most $\varepsilon$ on the same success
event. These leading constants are necessary
in the worst case even for arbitrary rank approximants.

\subsection{Classical inputs and organization}
The necessary matrix inequalities are classical. Araki supplies
log-majorization; Weyl--Horn and Kostant supply prescribed triangular
or Cholesky diagonals; successive Schur asymptotics are established
in matrix perturbation theory~\cite{Araki,Horn,LiMathias,Kostant,HT,UB,GAH}.
Appendix~\ref{ah:sec:araki-horn} gives their complete relative-range
combination, including signed pairs and batch envelopes. The main text
first gives a direct rank-one proof and then realizes its limits by
graphs. Series and parallel resistor laws, Schur complements,
Lanczos reduction, and Ritz min--max are classical; strongly coupled
graph coarse graining also has prior work~\cite{Koke}.

The degree-two proof uses Laptev's weighted-Cauchy spectral
asymptotic~\cite{Laptev}, a classical Fourier transform~\cite{DLMF},
and a uniform finite-grid comparison with one fixed continuum operator.
Stahl's rational approximation gives a separate nonsharp order
bound~\cite{Stahl}. The sharp higher-degree proof uses standard expander existence,
trace-class determinants, Kron reduction, and free/wired network
principles~\cite{Bordenave,Alon,Bornemann,GZ,DorflerBullo,LyonsPeres}. Rational
matrix-function updates are established~\cite{BCKS,SDA}. The algorithm
uses classical logarithmic quadrature and nearly linear graph
solvers~\cite{TW,KOSZ,KMP}. These credits do not certify priority for
the graph-relative specializations or their combination.

\input{sections/matrix-rank-one}
\input{sections/weighted-paths}
\input{sections/unweighted-insertion}
\input{sections/degree-dichotomy}
\input{sections/algorithm}
\input{sections/consequences}

\appendix
\input{sections/araki-horn}
\input{sections/finite-attainment}
\input{bibliography}
\end{document}
